\uuid{OJ8M}
\exo7id{4829}
\titre{exo7 4829}
\auteur{quercia}
\organisation{exo7}
\datecreate{2010-03-16}
\isIndication{false}
\isCorrection{true}
\chapitre{Topologie}
\sousChapitre{Compacité}
\cpge{mp-mpi.topologie-espaces-normes.g, mp-mpi.topologie-espaces-normes.e, pc.espaces-vectoriels-normes.e, psi.espaces-vectoriels-normes.e, mp-mpi.topologie-espaces-normes.d, pc.espaces-vectoriels-normes.d, psi.espaces-vectoriels-normes.d}
\module{Analyse}
\niveau{L2}
\difficulte{}

\contenu{
\question{
Soit $C$ un compact convexe d'un evn $E$. Soit~$f : C \to {C,}$
$1$-lipschitzienne.  Montrer que~$f$ admet un point fixe. On pourra
utiliser la fonction $f_n$~:~$x \mapsto\frac an +
\Bigl(1-\frac1n\Bigr)f(x)$ avec $a\in C$.
}
\reponse{
$C$ est stable par~$f_n$ qui est $\Bigl(1-\frac1n\Bigr)$-lipschitzienne.
Donc il existe $x_n\in C$ tel que $f_n(x_n) = x_n$~; toute valeur d'adh{\'e}rence de $(x_n)$
est point fixe de~$f$.
}
}
